\subsection{Scope and limitations}
  \label{subsec:discussion-scope-and-limitations}

  The present study is restricted to the emergence of geometric and
  gravitational structures.
  Matter degrees of freedom, energetic quantities, quantum statistics,
  and particle properties lie beyond the geometric regime addressed here
  and require additional layers of effective description.

  Configurations that do not admit a stable factorizable or projectable
  regime are likewise outside the scope of this analysis.
  In such cases, spacetime ceases to function as a meaningful
  descriptive language and must be replaced by a purely relational one.

  \paragraph{Kinematic status and constraints on dynamics.}
    The framework is deliberately kinematic: no equations of motion,
    Lagrangian, or Hamiltonian are postulated for the relational substrate.
    Its purpose is to identify the structural conditions under which an
    effective geometric description becomes admissible, independently of
    any specific dynamics.

    This does not render the framework non-predictive.
    The spectral invariants identified here — in particular the ratio
    $\lambda_2/\lambda_1 = 8/3$ in four-dimensional projectable regimes —
    impose necessary constraints on any consistent dynamical completion.
    Any admissible dynamics must reproduce these invariants within the
    geometric regime, while dynamics that generically violate spectral
    admissibility cannot be regarded as compatible extensions.

    As in general relativity, where geometry is specified prior to
    dynamical field equations, the present construction identifies the
    admissible geometric sector without prescribing the dynamics that
    selects it.
